\section{Non-response, dropout, relapse and harm}\label{sec:nonresponse}
\paragraph{Dropout} Dropout is the most visible failure. A meta-analysis of 115 CBT studies (20,995 participants) found 15.9\% dropout before treatment and 26.2\% during it, highest in depression and in e-therapy \citep{Fernandez2015Dropout}. CBT is not unusual here: across 669 psychotherapy studies the weighted dropout was 19.7\% and was not moderated by therapy orientation \citep{Swift2012Premature}; in 54 depression trials it was 17.5\% in psychotherapy arms \citep{Cooper2015DropoutMDD}, and in routine-care CBT for depression about 25\% \citep{Hans2013Depression}. CBT dropout is, if anything, lower than medication dropout \citep{Cuijpers2013Enduring}. Reported predictors are a weak alliance ($d=.55$; \citealp{Sharf2010Alliance}), comorbid personality disorder and minority status \citep{Cooper2015DropoutMDD}, and, in self-guided online programmes, male sex, low education and comorbid anxiety \citep{Karyotaki2015Dropout}.

\paragraph{Non-response} Even under trial conditions many patients do not respond. In a moderate-to-severe depression trial, 16-week response was 58\% and remission 40\% for cognitive therapy, with outcomes depending on therapist experience at one site \citep{DeRubeis2005CT}. Across 92 psychotherapy studies fewer than half of patients met response or remission criteria, and 48\% of care-as-usual patients also no longer met depression criteria, so the added value was about 14 percentage points \citep{Cuijpers2014Remission}. For anxiety, one review puts non-response at about a third of completers and finds that ``no predictor has been consistently supported'' \citep{Taylor2012NonResponse}; see also \citet{loerinc2015response} and \citet{springer2018remission}. Routine services do less well than trials: of 19,395 English IAPT patients, 40.3\% reliably recovered, 63.7\% reliably improved and 6.6\% reliably deteriorated \citep{Gyani2013IAPT}. The prognostic factors with support (comorbid personality disorder, chronicity, adversity) apply to all treatments (Section~\ref{sec:patients}); baseline severity did not moderate CBT versus medication in an individual-participant meta-analysis of 1,700 patients \citep{Weitz2015Severity}. In fairness, CBT-based collaborative care beat usual care in primary-care anxiety (12-month response 63.7\% versus 44.7\%; \citealp{RoyByrne2010CALM}) and adjunctive CBT roughly doubled response in treatment-resistant depression (46\% versus 22\%; \citealp{Wiles2013CoBalT}), with a benefit still present at 3--5 years \citep{Wiles2016CoBalTLong}; the Cochrane review rates the evidence for this as resting mostly on one trial \citep{Ijaz2018TRD}.

\paragraph{Relapse and durability} CBT has an enduring effect relative to stopping medication (relapse OR 2.61 favouring CBT; no significant difference from continued medication; \citealp{Cuijpers2013Enduring}; \citealp{Hollon2005Relapse}). Yet after acute-phase cognitive therapy ends, 29\% of patients relapse within one year and 54\% within two \citep{Vittengl2007Relapse}; continuation-phase CBT reduces this, especially in high-risk groups \citep{Jarrett2001Continuation}, and preventive psychological interventions cut relapse against usual care (RR 0.64; \citealp{Biesheuvel2015Recurrence}). In anxiety the follow-up record is thin: effects were not significant at 12 or more months for panic disorder, and relapse was reported in only 6 of 69 trials \citep{VanDis2020LongTerm}. In a panic trial CBT-alone response fell from 48.7\% acutely to 31.9\% six months after stopping \citep{Barlow2000Panic}.

\paragraph{Deterioration and harm} Across guided internet-CBT trials deterioration was lower with treatment than in controls (5.8\% versus 17.4\%, \citealp{Rozental2017Deterioration}; 3.36 versus 7.60, \citealp{Ebert2016Deterioration}), but low-education patients were at higher risk and did not differ from controls in the latter. Broader unwanted effects are common and poorly recorded: 5.2\% of 14,587 English therapy patients reported lasting bad effects \citep{Crawford2016Negative}; unwanted events are expected in roughly 5--20\% of patients \citep{Linden2014Rate}; therapists of 100 CBT patients reported side effects in 43 \citep{SchermulyHaupt2018Unwanted}; and a review of 85 trials found harmful events in more than 1 in 10 participants \citep{Klatte2023Adverse}. Only 21\% of 132 psychological-intervention trials monitored harms \citep{Jonsson2014Harms}, and none of 82 NIHR psychological trials mentioned an adverse event in the final report \citep{Duggan2014Adverse}. These are general to psychological therapy \citep{Lilienfeld2007Harm}, and the CBT literature has itself called for better harm detection \citep{Barlow2010Negative}. Routine outcome feedback for patients who are off track is a cheap mitigation (d = 0.23; \citealp{Delgadillo2018Feedback}; \citealp{Lambert2018Feedback}).

\paragraph{Reading} Dropout, non-response, relapse and harm are real, quantifiable and under-reported, so R4--R6 stand. They show that CBT leaves a large minority behind, not that it fails relative to alternatives, and almost none of the sources isolates the cognitive component.
